% Options for packages loaded elsewhere
\PassOptionsToPackage{unicode}{hyperref}
\PassOptionsToPackage{hyphens}{url}
\documentclass[
  a4paper,
]{ltjsarticle}
\usepackage{xcolor}
\usepackage[margin=25mm]{geometry}
\usepackage{amsmath,amssymb}
\setcounter{secnumdepth}{-\maxdimen} % remove section numbering
\usepackage{iftex}
\ifPDFTeX
  \usepackage[T1]{fontenc}
  \usepackage[utf8]{inputenc}
  \usepackage{textcomp} % provide euro and other symbols
\else % if luatex or xetex
  \usepackage{unicode-math} % this also loads fontspec
  \defaultfontfeatures{Scale=MatchLowercase}
  \defaultfontfeatures[\rmfamily]{Ligatures=TeX,Scale=1}
\fi
\usepackage{lmodern}
\ifPDFTeX\else
  % xetex/luatex font selection
\fi
% Use upquote if available, for straight quotes in verbatim environments
\IfFileExists{upquote.sty}{\usepackage{upquote}}{}
\IfFileExists{microtype.sty}{% use microtype if available
  \usepackage[]{microtype}
  \UseMicrotypeSet[protrusion]{basicmath} % disable protrusion for tt fonts
}{}
\makeatletter
\@ifundefined{KOMAClassName}{% if non-KOMA class
  \IfFileExists{parskip.sty}{%
    \usepackage{parskip}
  }{% else
    \setlength{\parindent}{0pt}
    \setlength{\parskip}{6pt plus 2pt minus 1pt}}
}{% if KOMA class
  \KOMAoptions{parskip=half}}
\makeatother
\usepackage{longtable,booktabs,array}
\newcounter{none} % for unnumbered tables
\usepackage{calc} % for calculating minipage widths
% Correct order of tables after \paragraph or \subparagraph
\usepackage{etoolbox}
\makeatletter
\patchcmd\longtable{\par}{\if@noskipsec\mbox{}\fi\par}{}{}
\makeatother
% Allow footnotes in longtable head/foot
\IfFileExists{footnotehyper.sty}{\usepackage{footnotehyper}}{\usepackage{footnote}}
\makesavenoteenv{longtable}
\usepackage{graphicx}
\makeatletter
\newsavebox\pandoc@box
\newcommand*\pandocbounded[1]{% scales image to fit in text height/width
  \sbox\pandoc@box{#1}%
  \Gscale@div\@tempa{\textheight}{\dimexpr\ht\pandoc@box+\dp\pandoc@box\relax}%
  \Gscale@div\@tempb{\linewidth}{\wd\pandoc@box}%
  \ifdim\@tempb\p@<\@tempa\p@\let\@tempa\@tempb\fi% select the smaller of both
  \ifdim\@tempa\p@<\p@\scalebox{\@tempa}{\usebox\pandoc@box}%
  \else\usebox{\pandoc@box}%
  \fi%
}
% Set default figure placement to htbp
\def\fps@figure{htbp}
\makeatother
\setlength{\emergencystretch}{3em} % prevent overfull lines
\providecommand{\tightlist}{%
  \setlength{\itemsep}{0pt}\setlength{\parskip}{0pt}}
\usepackage{bookmark}
\IfFileExists{xurl.sty}{\usepackage{xurl}}{} % add URL line breaks if available
\urlstyle{same}
\hypersetup{
  hidelinks,
  pdfcreator={LaTeX via pandoc}}

\author{}
\date{}

\begin{document}

\section{N体完全二体関係波における生成子ランクの線形上界と空間方向読出しの三方向飽和}\label{nux4f53ux5b8cux5168ux4e8cux4f53ux95a2ux4fc2ux6ce2ux306bux304aux3051ux308bux751fux6210ux5b50ux30e9ux30f3ux30afux306eux7ddaux5f62ux4e0aux754cux3068ux7a7aux9593ux65b9ux5411ux8aadux51faux3057ux306eux4e09ux65b9ux5411ux98fdux548c}

\subsection{関係波は二乗的に増え、回転モードは線形に増え、一意に読める空間方向は三で止まる}\label{ux95a2ux4fc2ux6ce2ux306fux4e8cux4e57ux7684ux306bux5897ux3048ux56deux8ee2ux30e2ux30fcux30c9ux306fux7ddaux5f62ux306bux5897ux3048ux4e00ux610fux306bux8aadux3081ux308bux7a7aux9593ux65b9ux5411ux306fux4e09ux3067ux6b62ux307eux308b}

\textbf{著者:} 木原 範昭\\
\textbf{ORCID:} 0009-0004-6753-4020\\
\textbf{日付:} 2026年7月21日\\
\textbf{Version DOI:} 10.5281/zenodo.21465899\\
\textbf{Concept DOI:} 10.5281/zenodo.21465898\\
\textbf{位置づけ:} 「次元の生成構造」シリーズ・第3論文 v1

\begin{center}\rule{0.5\linewidth}{0.5pt}\end{center}

\subsection{要旨}\label{ux8981ux65e8}

先行研究では、AB二体閉鎖位相系の一次元的な位相関係読出し［2,3］を、ABC三体・ABCD四体の完全二体関係波へ拡張し、ABCで一回転平面とZ不変法線、ABCDで三回転平面によるXYZ三方向の読出しが成立することを示した［4］。本稿はこの構成を一般の\(N\)体へ拡張し、内部関係波数・回転モード数・一意に読める空間方向数という三つの量が、それぞれ異なるオーダーで増えることを理論と数値実験で示す。

\(N\)体の完全二体関係波数は \(M_N=\binom{N}{2}=O(N^2)\)
で増える。一方、位相差正弦型の反対称生成子 \(\widetilde K\)
は、接続行列の行から作られるランク2以下の反対称行列\(N\)個の和

\[
\widetilde K=\sum_{k=1}^{N}\left(c_k s_k^{\mathsf T}-s_k c_k^{\mathsf T}\right)
\]

へ厳密に分解され、したがって

\[
\operatorname{rank}K\le 2\min\!\left(N,\left\lfloor\frac{M_N}{2}\right\rfloor\right)=O(N)
\]

が全\(N\)で成立する。さらに\(3\le N\le12\)については、tan半角パラメータによる厳密有理数の証人配置と実解析関数の零点集合の測度零性［9］を組み合わせ、ほとんど全ての初期位相で等号が成立することを証明した。最後に、二つの独立な関係方向が張る面から法線を読む方向読出しにおいて、法線候補空間は表示次元\(d\)に対し\(d-2\)次元であり、法線が符号を除いて一意に定まるのは\(d=3\)の場合に限ることを示す。したがって、現在の二体位相関係による読出し規則のもとでは、\(N\)体化で内部関係波と内部回転モードが増えても、一意に読める空間方向は三方向で飽和する。

数値実験では、\(N=3\)〜\(9\)の各32試行・720ステップで、二乗閉鎖誤差最大
\(1.82\times10^{-13}\)、絶対値二乗和変動最大
\(2.25\times10^{-13}\)、名称置換共変誤差最大 \(1.12\times10^{-13}\)
を確認し、頂点分解恒等式の残差は最大 \(5.00\times10^{-16}\)
だった。一般位置ランク則は\(N=3\)〜\(12\)の計2560試行で例外なく成立した。\(N\ge6\)で再出現する多次元零空間は、射影子と二乗量が一意である一方、内部基底が一意に選べない残余部分空間であることを直接確認した。状態の逐次正規化、観測減衰、絶対背景軸は用いていない。

\begin{center}\rule{0.5\linewidth}{0.5pt}\end{center}

\subsection{0. 結論}\label{ux7d50ux8ad6}

\(N\)体の完全二体関係波系では、次の三層が異なるオーダーで増える。

\[
\boxed{
M_N=\binom{N}{2}=O(N^2)
\quad\longrightarrow\quad
r\le N=O(N)
\quad\longrightarrow\quad
D_{\mathrm{unique}}=3=O(1)
}
\]

第一層は内部関係波数である。\(N\)体化で二乗的に増える。

第二層は内部回転モード数 \(r=\frac12\operatorname{rank}K\)
である。頂点分解定理により
\(\operatorname{rank}K\le2\min(N,\lfloor M_N/2\rfloor)\)
が全\(N\)で成立し、高々線形にしか増えない。\(3\le N\le12\)では一般位置でこの上界が等号で達成されることを厳密に証明した。

第三層は一意に読める空間方向数である。二関係面の法線候補空間は表示次元\(d\)に対し\(d-2\)次元であり、符号を除いた一意性は\(d=3\)の場合に限る。一般には
\(D_{\mathrm{unique}}\le3\)
であり、二つの独立かつ区別可能な関係方向が成立する一般位置で
\(D_{\mathrm{unique}}=3\) となる。\(N=2\)
や生成子が消滅する縮退位相では等号は成立しない。したがって、

\[
\boxed{
\text{現在の二体位相関係読出し規則のもとでは、}
\quad
\text{一意に読める空間方向は三方向で飽和する}
}
\]

\(N\ge6\)では零空間が再出現し、その次元は一般位置で

\[
\dim\ker K=M_N-2N=\frac{N(N-5)}{2}=O(N^2)
\]

として成長する。この零空間は、射影子と二乗量は一意に読めるが、内部の各方向を一意に選べない残余部分空間である。追加関係波は消えるのではなく、この基底を持たない内部部分空間と、空間方向として採用されない内部位相モードとして保持される。

これらは基本公理系v5の公理16が定める採用条件、すなわち回転周波数による区別可能性と面法線の一意性の実装である［1］。

\begin{center}\rule{0.5\linewidth}{0.5pt}\end{center}

\subsection{1. 研究課題}\label{ux7814ux7a76ux8ab2ux984c}

\subsubsection{1.1
先行論文から残った問い}\label{ux5148ux884cux8ad6ux6587ux304bux3089ux6b8bux3063ux305fux554fux3044}

先行するABCD四体実験［4］は、六つの完全二体関係波が三つの回転平面へ分解し、XYZ三方向と読める構造が生じることを示した。同時に、次の問いが残った。

\begin{quote}
五体以上へ関係波を増やしたとき、内部で増える量は何であり、空間方向として一意に読める量はどこで止まるのか。
\end{quote}

［4］は「五体以上で増えるのは方向を一意に定められない内部関係である」という解釈を示したが、これは数値実験の対象外であり、定理としては示されていなかった。本稿はこの解釈を、頂点分解によるランク上界、一般位置での等号、法線一意性の三つの定理へ昇格させる。

\subsubsection{1.2
因果方向の規則}\label{ux56e0ux679cux65b9ux5411ux306eux898fux5247}

本稿の因果方向は次である。

\[
\boxed{
\text{関係波と二乗閉鎖}
\longrightarrow
\text{読出し写像}
\longrightarrow
\text{一意に再構成できる方向}
\longrightarrow
\text{空間表示}
}
\]

空間は最初に存在しない。したがって本稿では次を行わない。

\begin{enumerate}
\def\labelenumi{\arabic{enumi}.}
\tightlist
\item
  関係波 \(X_{ij}\) を背景空間の距離 \(d_{ij}\) と読み替えない。
\item
  既存Euclidean空間への埋め込みランクを、空間生成の判定に使わない。
\item
  回転平面数、生成子ランク、零空間次元を、空間軸数と同一視しない。
\item
  五体目以降に外部幾何由来の拘束を追加しない。全関係波は同じ一つの二乗閉鎖へ参加する。
\end{enumerate}

本稿が判定するのは、状態空間の大きさではなく、読出し写像の一意性である。

\subsubsection{1.3
本稿で判別すること}\label{ux672cux7a3fux3067ux5224ux5225ux3059ux308bux3053ux3068}

\begin{enumerate}
\def\labelenumi{\arabic{enumi}.}
\tightlist
\item
  生成子ランクの全\(N\)上界を、モデルの内部から導出できるか。
\item
  一般位置でその上界が等号になるか。厳密に証明できるか。
\item
  法線読出しの一意性は、どの表示次元で成立するか。
\item
  \(N\ge6\)で再出現する零空間は、どのような量として読めるか。
\item
  保存量と名称置換共変性は\(N\)拡張後も維持されるか。
\end{enumerate}

\begin{center}\rule{0.5\linewidth}{0.5pt}\end{center}

\subsection{2. 主張の分類}\label{ux4e3bux5f35ux306eux5206ux985e}

本稿では、定義、数学的帰結、数値実験事実、空間方向解釈、未固定事項を区別する。無名性・二乗閉鎖・完全二体関係波の公理的基盤は基本公理系v5に基づく［1］。

{\def\LTcaptype{none} % do not increment counter
\begin{longtable}[]{@{}
  >{\raggedright\arraybackslash}p{(\linewidth - 4\tabcolsep) * \real{0.3333}}
  >{\raggedright\arraybackslash}p{(\linewidth - 4\tabcolsep) * \real{0.3333}}
  >{\raggedright\arraybackslash}p{(\linewidth - 4\tabcolsep) * \real{0.3333}}@{}}
\toprule\noalign{}
\begin{minipage}[b]{\linewidth}\raggedright
対象
\end{minipage} & \begin{minipage}[b]{\linewidth}\raggedright
分類
\end{minipage} & \begin{minipage}[b]{\linewidth}\raggedright
本稿での位置づけ
\end{minipage} \\
\midrule\noalign{}
\endhead
\bottomrule\noalign{}
\endlastfoot
完全二体関係波を独立な状態波とする & モデル定義 &
［1］公理16の実装、［4］の継承 \\
\(\sum_e X_e^2+(iR)^2=0\) & 閉鎖条件 & 各ステップで直接計算 \\
端点共有関係だけを結合する位相差正弦生成子 & 作業仮説 &
［4］と同一。公理からは未導出 \\
頂点分解 \(\widetilde K=\sum_k(c_ks_k^{\mathsf T}-s_kc_k^{\mathsf T})\)
& 数学的帰結 & 定理1。数値残差 \(5.00\times10^{-16}\) \\
\(\operatorname{rank}K\le2\min(N,\lfloor M/2\rfloor)\)（全\(N\)） &
数学的帰結 & 定理1の系 \\
一般位置での等号（\(3\le N\le12\)） & 数学的帰結（計算機援用証明） &
定理2。厳密有理数証人＋測度零論法 \\
一般位置での等号（全\(N\)） & 新仮説 & 一般証人構成は未証明 \\
法線一意性（\(d=3\)に限る） & 数学的帰結 & 定理3 \\
三方向飽和（\(D_{\mathrm{unique}}\le3\)、一般位置で\(=3\)） &
公理16と定理3からの帰結 & 第7節。生成子ランク単独からは出ない \\
実験A〜Eの結果 & 数値実験事実 & 第8節 \\
三区別可能回転平面をXYZと読む & 空間方向解釈 & ［4］の解釈の継承 \\
残余方向を内部位相モード・残余部分空間と読む & 空間方向解釈 &
［1］公理16の残余規定と整合 \\
残余方向が対応する物理軸 & 対象外 & 本稿では同定しない \\
\end{longtable}
}

\begin{center}\rule{0.5\linewidth}{0.5pt}\end{center}

\subsection{3.
N体完全二体関係波モデル}\label{nux4f53ux5b8cux5168ux4e8cux4f53ux95a2ux4fc2ux6ce2ux30e2ux30c7ux30eb}

\subsubsection{3.1
関係集合と関係波}\label{ux95a2ux4fc2ux96c6ux5408ux3068ux95a2ux4fc2ux6ce2}

\(N\)個の無名な個体に対し、全二体関係の集合を

\[
\mathcal E_N
=
\bigl\{\{i,j\}\mid 1\le i<j\le N\bigr\},
\qquad
M=\lvert\mathcal E_N\rvert=\binom{N}{2}
\]

とする。各関係 \(e\) に複素関係波 \(X_e=a_e+ib_e\) を対応させ、全状態を
\(X\in\mathbb C^M\)
と書く。関係波は観測チャネルではなく、成分波から導出されない独立な物理状態成分である［1］。

閉鎖条件は

\[
\sum_{e\in\mathcal E_N}X_e^2+(iR)^2=0,
\qquad
\text{移項表示}\quad
q(X)=X^{\mathsf T}X=R^2
\]

である。比較用の独立な保存量として絶対値二乗和 \(H(X)=X^\dagger X\)
も記録する。

\subsubsection{3.2
隣接と生成子}\label{ux96a3ux63a5ux3068ux751fux6210ux5b50}

二関係 \(e,f\) が一つの端点を共有するときだけ隣接とする。

\[
A_{ef}
=
\begin{cases}
1,& e\ne f\ \text{かつ}\ e\cap f\ne\varnothing,\\
0,& \text{それ以外}.
\end{cases}
\]

これは完全グラフ \(K_N\) の線グラフの隣接行列である［5］。

初期位相 \(\theta_e=\arg X_e(0)\) から、生成子の未規格化形を

\[
\widetilde K_{ef}
=
A_{ef}\sin(\theta_f-\theta_e)
\]

と置く。\(\widetilde K^{\mathsf T}=-\widetilde K\)
であり、非零ならスペクトルノルムで割って \(K\)
とする。この位相差正弦結合則は、先行論文［4］と同一の保存的な作業仮説である。本稿の動力学は固定生成子である。生成子は初期位相から一回構成し、以後固定する。

一ステップ更新はCayley変換

\[
U=(I-\gamma K)^{-1}(I+\gamma K),
\qquad
U^{\mathsf T}U=I
\]

による実直交行列とする［7］。直交性により \(q(X)\) と \(H(X)\)
は逐次正規化なしで保存される。

\subsubsection{3.3 接続行列}\label{ux63a5ux7d9aux884cux5217}

個体を行、関係を列とする符号なし接続行列を \(B\in\mathbb R^{N\times M}\)
とする。関係 \(e=\{i,j\}\) の列は、行 \(i,j\)
だけが1である。線グラフ隣接行列とは

\[
A=B^{\mathsf T}B-2I_M
\]

の関係にある。

\begin{center}\rule{0.5\linewidth}{0.5pt}\end{center}

\subsection{4. 頂点分解定理}\label{ux9802ux70b9ux5206ux89e3ux5b9aux7406}

\subsubsection{4.1 定理}\label{ux5b9aux7406}

\begin{quote}
\textbf{定理1（頂点分解）}

\(B\)の第\(k\)行を\(b_k\)とし、

\[
c_k=(B_{ke}\cos\theta_e)_{e\in\mathcal E_N},
\qquad
s_k=(B_{ke}\sin\theta_e)_{e\in\mathcal E_N}
\]

と置く。このとき

\[
\widetilde K=\sum_{k=1}^{N}\left(c_k s_k^{\mathsf T}-s_k c_k^{\mathsf T}\right)
\]

が厳密に成立する。
\end{quote}

\textbf{証明} 成分計算による。

\[
\left(c_k s_k^{\mathsf T}-s_k c_k^{\mathsf T}\right)_{ef}
=
B_{ke}B_{kf}\left(\cos\theta_e\sin\theta_f-\sin\theta_e\cos\theta_f\right)
=
B_{ke}B_{kf}\sin(\theta_f-\theta_e)
\]

\(k\)について和を取ると、係数は
\(\sum_kB_{ke}B_{kf}=(B^{\mathsf T}B)_{ef}\) となる。\(e=f\)では
\(\sin0=0\)
で対角は消える。\(e\ne f\)では、相異なる二辺が共有できる端点は高々一つなので
\((B^{\mathsf T}B)_{ef}=A_{ef}\in\{0,1\}\)。したがって全成分が
\(A_{ef}\sin(\theta_f-\theta_e)=\widetilde K_{ef}\) に一致する。∎

数値実験では、この恒等式の残差は全構成・全試行で最大
\(5.00\times10^{-16}\)（機械精度）だった。

\subsubsection{4.2
ランクの線形上界}\label{ux30e9ux30f3ux30afux306eux7ddaux5f62ux4e0aux754c}

\begin{quote}
\textbf{系1（ランク上界）}

全ての\(N\)と全ての初期位相に対し、

\[
\operatorname{rank}K
\le
2\min\!\left(N,\left\lfloor\frac{M}{2}\right\rfloor\right)
\]
\end{quote}

\textbf{証明} 定理1の各項 \(c_ks_k^{\mathsf T}-s_kc_k^{\mathsf T}\)
はランク2以下なので、和のランクは \(2N\)
以下である。\(\operatorname{rank}K\le M\)
は自明である。実反対称行列は直交変換で二次元回転ブロックの直和と零ブロックへ分解されるため、ランクは偶数である［6］。三つを合わせて上界を得る。∎

系1は、関係波数が \(O(N^2)\) で増えるのに対し、動的な回転自由度が高々
\(O(N)\) であることを、背景空間を導入せずモデルの内部だけから示す。

\textbf{補足（上界が飽和可能である構造的理由）} \(K_N\) は \(N\ge3\)
で連結かつ非二部であり、符号なし接続行列の実数体上のランクは \(N\)
である［8］。したがって \(\widetilde K\) の列空間を張る \(2N\)
本のベクトル \(\{c_k,s_k\}\) は、\(N\)
本の独立な接続行から作られており、係数空間の側に退化がない。次節で示すとおり、一般位置ではこの上界が実際に達成される。

\begin{center}\rule{0.5\linewidth}{0.5pt}\end{center}

\subsection{5.
一般位置ランク定理}\label{ux4e00ux822cux4f4dux7f6eux30e9ux30f3ux30afux5b9aux7406}

\subsubsection{5.1 定理}\label{ux5b9aux7406-1}

\begin{quote}
\textbf{定理2（一般位置等号、\(3\le N\le12\)）}

\(3\le N\le12\)の各\(N\)に対し、\(m=\min(N,\lfloor M/2\rfloor)\)と置くと、ほとんど全ての（Lebesgue測度零の例外集合を除く全ての）初期位相
\(\theta\in[0,2\pi)^M\) に対し、

\[
\operatorname{rank}K=2m
\]
\end{quote}

\textbf{証明} 上界は系1である。下界を示す。\(\widetilde K\)の成分は

\[
\sin(\theta_f-\theta_e)
=
\sin\theta_f\cos\theta_e-\cos\theta_f\sin\theta_e
\]

であり、単位円上の点 \((\cos\theta_e,\sin\theta_e)\)
の多項式である。tan半角パラメータ \(t_e\) により

\[
(\cos\theta_e,\sin\theta_e)
=
\left(\frac{1-t_e^2}{1+t_e^2},\ \frac{2t_e}{1+t_e^2}\right)
\]

と有理パラメータ化すると、\(t_e=e\)（\(e=1,\dots,M\)）の証人配置で
\(\widetilde K\)
は厳密な有理行列になる。この行列のランクを分数演算のガウス消去で厳密に計算した結果、\(3\le N\le12\)の全てで

\[
\operatorname{rank}\widetilde K=2m
\]

が確認された（第10節の検証プログラムと出力）。浮動小数を含まない厳密有理数演算であるから、これは各\(N\)に対する証人の存在証明である。

証人配置で \(\operatorname{rank}\widetilde K=2m\) であることから、ある
\(2m\times2m\) 小行列式 \(D(\theta)\) が証人で非零である。\(D\) は
\(\theta\)
の実解析関数であり、恒等的に零ではない。非自明な実解析関数の零点集合はLebesgue測度零である［9］。ランクが
\(2m\) 未満となる位相集合は \(\{\theta\mid D(\theta)=0\}\)
に含まれるため、測度零である。∎

\subsubsection{5.2
全Nへの一般化}\label{ux5168nux3078ux306eux4e00ux822cux5316}

\begin{quote}
\textbf{新仮説（一般位置等号、全\(N\)）}

全ての\(N\ge3\)に対し、一般位置で
\(\operatorname{rank}K=2\min(N,\lfloor M/2\rfloor)\) が成立する。
\end{quote}

\(N\le12\)では定理2で証明済みである。任意の\(N\)に対する証人の一般構成は未証明であり、本稿では新仮説として明示する。同じ厳密有理数検証は任意の個別\(N\)へ機械的に拡張できる。

\subsubsection{5.3 縮退配置}\label{ux7e2eux9000ux914dux7f6e}

一般位置から外れる配置では、ランクは低下し得る。数値実験（実験B）では次を確認した。

\begin{enumerate}
\def\labelenumi{\arabic{enumi}.}
\tightlist
\item
  全関係波が等位相、または位相が \(\{0,\pi\}\)
  に限られる場合、\(\sin(\theta_f-\theta_e)=0\)
  となり生成子は消滅する（ランク0）。
\item
  位相が二値 \(\{0,0.5\}\) の場合、\(N=5\) でのみランクが \(10\) から
  \(8\) へ低下した。\(N=5\) は \(M=2N\)
  の臨界構成であり、この観察は定理2の例外集合が空でないことの具体例である。
\end{enumerate}

\begin{center}\rule{0.5\linewidth}{0.5pt}\end{center}

\subsection{6.
法線一意性定理}\label{ux6cd5ux7ddaux4e00ux610fux6027ux5b9aux7406}

\subsubsection{6.1
一意な方向の定義}\label{ux4e00ux610fux306aux65b9ux5411ux306eux5b9aux7fa9}

方向が一意に定まるとは、次を満たすことをいう。

\begin{enumerate}
\def\labelenumi{\arabic{enumi}.}
\tightlist
\item
  方向が関係データだけから再構成できる。
\item
  絶対背景軸を参照しない。
\item
  個体名の置換に対して共変に移る。
\item
  任意の候補から一本を選ぶ追加規則を必要としない。
\end{enumerate}

ここでの一意性は、まず符号を除いた一次元部分空間の一意性である。

\subsubsection{6.2
二関係面と法線候補空間}\label{ux4e8cux95a2ux4fc2ux9762ux3068ux6cd5ux7ddaux5019ux88dcux7a7aux9593}

二つの一次独立な関係方向
\(u,v\in V_d\)（\(V_d\)は実\(d\)次元表示空間）が張る面を
\(P=\operatorname{span}\{u,v\}\) とする。行列
\(A_P=\begin{pmatrix}u^{\mathsf T}\\ v^{\mathsf T}\end{pmatrix}\)
に対し、法線候補空間は

\[
\mathcal N(P)=\ker A_P=P^\perp
\]

であり、その直交射影

\[
P_\perp=I_d-A_P^{\mathsf T}\left(A_PA_P^{\mathsf T}\right)^{-1}A_P
\]

は外部の軸名を使わず \(u,v\) だけから定まる。

\subsubsection{6.3 定理}\label{ux5b9aux7406-2}

\begin{quote}
\textbf{定理3（法線一意性）}

二つの一次独立な関係方向が張る面から、追加の背景軸または選択規則なしに法線方向を再構成する場合、法線が符号を除いて一意に定まるのは、表示次元が

\[
d=3
\]

の場合に限る。
\end{quote}

\textbf{証明} \(\operatorname{rank}A_P=2\)
なので、階数・零化次数定理により \(\dim\mathcal N(P)=d-2\) である。

\(d=3\)では \(\dim\mathcal N(P)=1\) であり、規格化法線 \(n\)
による射影子 \(P_\perp=nn^{\mathsf T}\) は
\((-n)(-n)^{\mathsf T}=nn^{\mathsf T}\)
で符号反転に不変である。よって方向直線は一意である。

\(d\ge4\)では \(\dim\mathcal N(P)=d-2\ge2\)
である。\(u,v\)を固定する直交変換の群は法線候補空間上で \(O(d-2)\)
として作用し、\(P_\perp\)
を変えないまま法線候補を互いに混合する。単位法線候補の集合は
\(S^{d-3}\)、符号同一視後も \(\mathbb{RP}^{d-3}\)
という連続族であり、\(d=3\)の \(\mathbb{RP}^0\)（一点）に対し
\(d\ge4\)では一点でない。関係データはこの族の要素を区別しないため、一本を選ぶことは無名性と両立しない。∎

\(d\ge4\)でも、射影子 \(P_\perp\)
と候補成分の二乗和は一意に読める。一意に読めないのは、その内部の各一次方向である。

\begin{center}\rule{0.5\linewidth}{0.5pt}\end{center}

\subsection{7.
三層階層と残余部分空間}\label{ux4e09ux5c64ux968eux5c64ux3068ux6b8bux4f59ux90e8ux5206ux7a7aux9593}

\subsubsection{7.1 三層の合成}\label{ux4e09ux5c64ux306eux5408ux6210}

定理1・定理2による内部モード構造と、公理16に基づく観測選択［1］へ定理3を適用すると、\(N\)体完全二体関係波系は次の三層に分離する。

{\def\LTcaptype{none} % do not increment counter
\begin{longtable}[]{@{}
  >{\raggedright\arraybackslash}p{(\linewidth - 6\tabcolsep) * \real{0.2500}}
  >{\raggedright\arraybackslash}p{(\linewidth - 6\tabcolsep) * \real{0.2500}}
  >{\raggedright\arraybackslash}p{(\linewidth - 6\tabcolsep) * \real{0.2500}}
  >{\raggedright\arraybackslash}p{(\linewidth - 6\tabcolsep) * \real{0.2500}}@{}}
\toprule\noalign{}
\begin{minipage}[b]{\linewidth}\raggedright
層
\end{minipage} & \begin{minipage}[b]{\linewidth}\raggedright
量
\end{minipage} & \begin{minipage}[b]{\linewidth}\raggedright
オーダー
\end{minipage} & \begin{minipage}[b]{\linewidth}\raggedright
根拠
\end{minipage} \\
\midrule\noalign{}
\endhead
\bottomrule\noalign{}
\endlastfoot
内部関係波 & \(M_N=\binom{N}{2}\) & \(O(N^2)\) & モデル定義 \\
内部回転モード & \(r=\frac12\operatorname{rank}K\le N\) & \(O(N)\) &
定理1・定理2 \\
一意に読める空間方向 & \(D_{\mathrm{unique}}\le3\)（一般位置で \(=3\)）
& \(O(1)\) & 公理16＋定理3 \\
\end{longtable}
}

三方向飽和は、生成子ランクだけの帰結ではない。公理16の観測選択（区別可能性＋法線一意性）へ定理3を接続して初めて成立する。また
\(D_{\mathrm{unique}}=3\)
は常に成立するのではない。二つの独立かつ区別可能な関係方向が成立する一般位置で
\(D_{\mathrm{unique}}=3\) となり、\(N=2\)
や生成子が消滅する縮退位相では等号は成立しない。

回転平面数は空間方向数ではない。回転平面は、独立な回転周波数によって区別できる場合に方向の候補となるが、空間方向としての採用は法線一意性を満たす三方向までに限られる。区別可能な回転平面が三を超えて増える場合、超える平面は内部位相モードとして保持される。これは基本公理系v5の公理16が定める採用条件の実装である［1］。

\subsubsection{7.2
残余部分空間の成長}\label{ux6b8bux4f59ux90e8ux5206ux7a7aux9593ux306eux6210ux9577}

一般位置では \(N\ge5\) で

\[
\dim\ker K
=
M_N-2N
=
\frac{N(N-5)}{2}
\]

となる（\(N=3\)では1、\(N=4,5\)では0）。\(N\ge6\)で零空間が再出現し、\(O(N^2)\)
で成長する。大きな\(N\)では、関係波空間の大部分が生成子に動かされない残余部分空間になる。

ABCの一次元核はZ法線として一意に読めた［4］。\(N\ge6\)の核は多次元であり、核への直交射影子と射影ノルムは一意に定まるが、核内部の各方向は
\(O(\dim\ker K)\)
のゲージ自由度を持ち、一意に選べない。したがって残余は、

\[
\boxed{
\text{一意な線形方向ではなく、}
\quad
\text{一意な基底を持たない内部部分空間として残る}
}
\]

と読む。これも公理16の残余規定と整合する［1］。

\subsubsection{7.3
何が保存されるか}\label{ux4f55ux304cux4fddux5b58ux3055ux308cux308bux304b}

\(N\)体化で維持されるのは、特定の名称を持つX、Y、Zではない。維持されるのは、区別可能な回転平面から選ばれる二つの独立な関係方向と、その関係面から一意に定まる一つの法線方向という三方向構造である。個体名を置換すれば三方向を構成する関係成分も共変に移る。

\begin{center}\rule{0.5\linewidth}{0.5pt}\end{center}

\subsection{8. 数値実験}\label{ux6570ux5024ux5b9fux9a13}

\subsubsection{8.1 共通条件}\label{ux5171ux901aux6761ux4ef6}

{\def\LTcaptype{none} % do not increment counter
\begin{longtable}[]{@{}lr@{}}
\toprule\noalign{}
項目 & 値 \\
\midrule\noalign{}
\endhead
\bottomrule\noalign{}
\endlastfoot
動力学 & 固定生成子（初期位相から一回構成） \\
\(R^2\) & 1.0 \\
虚部初期振幅 \(s\) & 0.35 \\
Cayley係数 \(\gamma\) & \(\tan(\pi/144)\) \\
乱数シード & 20260721（実験B〜Eは +1〜+4） \\
判定許容差 & \(10^{-10}\) \\
状態の逐次正規化 & なし \\
観測減衰 & なし \\
絶対背景軸 & なし \\
\end{longtable}
}

初期状態は、実単位ベクトル \(u\perp v\) を用いた
\(X(0)=\sqrt{R^2+s^2}\,u+isv\) であり、\(q(X(0))=R^2\)
が厳密に成立する［4］。

\subsubsection{8.2 実験A:
N=3〜9の保存量・共変性・スペクトル構造}\label{ux5b9fux9a13a-n39ux306eux4fddux5b58ux91cfux5171ux5909ux6027ux30b9ux30daux30afux30c8ux30ebux69cbux9020}

各\(N\)で32試行、各720ステップ。

{\def\LTcaptype{none} % do not increment counter
\begin{longtable}[]{@{}rrrrrl@{}}
\toprule\noalign{}
\(N\) & \(M\) & rank \(K\) & \(\dim\ker K\) & 回転平面数 & ランク則 \\
\midrule\noalign{}
\endhead
\bottomrule\noalign{}
\endlastfoot
3 & 3 & 2 & 1 & 1 & 成立 \\
4 & 6 & 6 & 0 & 3 & 成立 \\
5 & 10 & 10 & 0 & 5 & 成立 \\
6 & 15 & 12 & 3 & 6 & 成立 \\
7 & 21 & 14 & 7 & 7 & 成立 \\
8 & 28 & 16 & 12 & 8 & 成立 \\
9 & 36 & 18 & 18 & 9 & 成立 \\
\end{longtable}
}

ランクと零空間次元は各\(N\)の全32試行で一致し、全て
\(2\min(N,\lfloor M/2\rfloor)\) と \(M-2N\)（\(N\ge5\)）に一致した。

{\def\LTcaptype{none} % do not increment counter
\begin{longtable}[]{@{}
  >{\raggedleft\arraybackslash}p{(\linewidth - 12\tabcolsep) * \real{0.1429}}
  >{\raggedleft\arraybackslash}p{(\linewidth - 12\tabcolsep) * \real{0.1429}}
  >{\raggedleft\arraybackslash}p{(\linewidth - 12\tabcolsep) * \real{0.1429}}
  >{\raggedleft\arraybackslash}p{(\linewidth - 12\tabcolsep) * \real{0.1429}}
  >{\raggedleft\arraybackslash}p{(\linewidth - 12\tabcolsep) * \real{0.1429}}
  >{\raggedleft\arraybackslash}p{(\linewidth - 12\tabcolsep) * \real{0.1429}}
  >{\raggedleft\arraybackslash}p{(\linewidth - 12\tabcolsep) * \real{0.1429}}@{}}
\toprule\noalign{}
\begin{minipage}[b]{\linewidth}\raggedleft
\(N\)
\end{minipage} & \begin{minipage}[b]{\linewidth}\raggedleft
最大二乗閉鎖誤差
\end{minipage} & \begin{minipage}[b]{\linewidth}\raggedleft
最大絶対値二乗和変動
\end{minipage} & \begin{minipage}[b]{\linewidth}\raggedleft
最大直交誤差
\end{minipage} & \begin{minipage}[b]{\linewidth}\raggedleft
最大軌道置換共変誤差
\end{minipage} & \begin{minipage}[b]{\linewidth}\raggedleft
頂点分解残差
\end{minipage} & \begin{minipage}[b]{\linewidth}\raggedleft
最大核射影変動
\end{minipage} \\
\midrule\noalign{}
\endhead
\bottomrule\noalign{}
\endlastfoot
3 & \(1.44\times10^{-13}\) & \(1.78\times10^{-13}\) &
\(4.44\times10^{-16}\) & \(7.62\times10^{-14}\) & \(4.44\times10^{-16}\)
& \(2.81\times10^{-14}\) \\
4 & \(1.42\times10^{-13}\) & \(1.74\times10^{-13}\) &
\(6.66\times10^{-16}\) & \(1.12\times10^{-13}\) & \(4.44\times10^{-16}\)
& \(0\) \\
5 & \(1.82\times10^{-13}\) & \(2.25\times10^{-13}\) &
\(6.66\times10^{-16}\) & \(6.58\times10^{-14}\) & \(4.72\times10^{-16}\)
& \(0\) \\
6 & \(8.91\times10^{-14}\) & \(1.10\times10^{-13}\) &
\(6.66\times10^{-16}\) & \(4.36\times10^{-14}\) & \(5.00\times10^{-16}\)
& \(1.80\times10^{-14}\) \\
7 & \(1.28\times10^{-13}\) & \(1.58\times10^{-13}\) &
\(8.88\times10^{-16}\) & \(4.91\times10^{-14}\) & \(5.00\times10^{-16}\)
& \(2.81\times10^{-14}\) \\
8 & \(1.27\times10^{-13}\) & \(1.55\times10^{-13}\) &
\(1.33\times10^{-15}\) & \(6.42\times10^{-14}\) & \(5.00\times10^{-16}\)
& \(4.60\times10^{-14}\) \\
9 & \(1.70\times10^{-13}\) & \(2.21\times10^{-13}\) &
\(1.67\times10^{-15}\) & \(8.99\times10^{-14}\) & \(4.72\times10^{-16}\)
& \(4.44\times10^{-14}\) \\
\end{longtable}
}

全構成・全試行で許容差 \(10^{-10}\)
を下回った。生成子の名称置換共変誤差は最大 \(3.33\times10^{-16}\)
だった。零空間成分は多次元核（\(N\ge6\)）でも軌道全体で保存された。

\begin{figure}
\centering
\pandocbounded{\includegraphics[keepaspectratio,alt={保存量と共変性}]{figs/nbody_conservation_errors_v1.png}}
\caption{保存量と共変性}
\end{figure}

\subsubsection{8.3 実験B:
一般位置ランク則と縮退探索}\label{ux5b9fux9a13b-ux4e00ux822cux4f4dux7f6eux30e9ux30f3ux30afux5247ux3068ux7e2eux9000ux63a2ux7d22}

\(N=3\)〜\(12\)の各256試行（一様乱数位相、計2560試行）の全てで

\[
\operatorname{rank}K=2\min\!\left(N,\left\lfloor\frac{M}{2}\right\rfloor\right)
\]

が成立した。

{\def\LTcaptype{none} % do not increment counter
\begin{longtable}[]{@{}
  >{\raggedleft\arraybackslash}p{(\linewidth - 12\tabcolsep) * \real{0.1429}}
  >{\raggedleft\arraybackslash}p{(\linewidth - 12\tabcolsep) * \real{0.1429}}
  >{\raggedleft\arraybackslash}p{(\linewidth - 12\tabcolsep) * \real{0.1429}}
  >{\raggedleft\arraybackslash}p{(\linewidth - 12\tabcolsep) * \real{0.1429}}
  >{\raggedleft\arraybackslash}p{(\linewidth - 12\tabcolsep) * \real{0.1429}}
  >{\raggedleft\arraybackslash}p{(\linewidth - 12\tabcolsep) * \real{0.1429}}
  >{\raggedleft\arraybackslash}p{(\linewidth - 12\tabcolsep) * \real{0.1429}}@{}}
\toprule\noalign{}
\begin{minipage}[b]{\linewidth}\raggedleft
\(N\)
\end{minipage} & \begin{minipage}[b]{\linewidth}\raggedleft
\(M\)
\end{minipage} & \begin{minipage}[b]{\linewidth}\raggedleft
期待ランク
\end{minipage} & \begin{minipage}[b]{\linewidth}\raggedleft
一般位置観測
\end{minipage} & \begin{minipage}[b]{\linewidth}\raggedleft
等位相
\end{minipage} & \begin{minipage}[b]{\linewidth}\raggedleft
\(\{0,\pi\}\)位相
\end{minipage} & \begin{minipage}[b]{\linewidth}\raggedleft
二値\(\{0,0.5\}\)位相
\end{minipage} \\
\midrule\noalign{}
\endhead
\bottomrule\noalign{}
\endlastfoot
3 & 3 & 2 & 2 & 0 & 0 & 2 \\
4 & 6 & 6 & 6 & 0 & 0 & 6 \\
5 & 10 & 10 & 10 & 0 & 0 & \textbf{8} \\
6 & 15 & 12 & 12 & 0 & 0 & 12 \\
7 & 21 & 14 & 14 & 0 & 0 & 14 \\
8 & 28 & 16 & 16 & 0 & 0 & 16 \\
9 & 36 & 18 & 18 & 0 & 0 & 18 \\
10 & 45 & 20 & 20 & 0 & 0 & 20 \\
11 & 55 & 22 & 22 & 0 & 0 & 22 \\
12 & 66 & 24 & 24 & 0 & 0 & 24 \\
\end{longtable}
}

\begin{figure}
\centering
\pandocbounded{\includegraphics[keepaspectratio,alt={三層スケーリング}]{figs/three_layer_scaling_v1.png}}
\caption{三層スケーリング}
\end{figure}

\subsubsection{8.4
厳密有理数証人}\label{ux53b3ux5bc6ux6709ux7406ux6570ux8a3cux4eba}

定理2の証人配置 \(t_e=e\) に対し、分数演算のガウス消去で
\(\operatorname{rank}\widetilde K\)
を厳密に計算した。\(N=3\)〜\(12\)の全てで
\(\operatorname{rank}\widetilde K=2\min(N,\lfloor M/2\rfloor)\)
が厳密に成立した。この計算に浮動小数点は含まれない。

\subsubsection{8.5 実験C:
法線非一意性の実証}\label{ux5b9fux9a13c-ux6cd5ux7ddaux975eux4e00ux610fux6027ux306eux5b9fux8a3c}

表示次元 \(d=3,4,5\) の各32試行で、二つの一次独立方向 \(u,v\)
から法線候補空間を構成し、候補空間内の直交変換（ゲージ変換）64標本に対し、射影子の変動と候補法線間の最大角を測定した。

{\def\LTcaptype{none} % do not increment counter
\begin{longtable}[]{@{}
  >{\raggedleft\arraybackslash}p{(\linewidth - 8\tabcolsep) * \real{0.2105}}
  >{\raggedleft\arraybackslash}p{(\linewidth - 8\tabcolsep) * \real{0.2105}}
  >{\raggedleft\arraybackslash}p{(\linewidth - 8\tabcolsep) * \real{0.2105}}
  >{\raggedleft\arraybackslash}p{(\linewidth - 8\tabcolsep) * \real{0.2105}}
  >{\raggedright\arraybackslash}p{(\linewidth - 8\tabcolsep) * \real{0.1579}}@{}}
\toprule\noalign{}
\begin{minipage}[b]{\linewidth}\raggedleft
\(d\)
\end{minipage} & \begin{minipage}[b]{\linewidth}\raggedleft
候補空間次元 \(d-2\)
\end{minipage} & \begin{minipage}[b]{\linewidth}\raggedleft
射影子ゲージ変動
\end{minipage} & \begin{minipage}[b]{\linewidth}\raggedleft
候補法線間最大角
\end{minipage} & \begin{minipage}[b]{\linewidth}\raggedright
符号を除き一意
\end{minipage} \\
\midrule\noalign{}
\endhead
\bottomrule\noalign{}
\endlastfoot
3 & 1 & \(4.44\times10^{-16}\) & \(0.0000^\circ\) & 成立 \\
4 & 2 & \(9.99\times10^{-16}\) & \(89.9986^\circ\) & 不成立 \\
5 & 3 & \(1.11\times10^{-15}\) & \(89.9998^\circ\) & 不成立 \\
\end{longtable}
}

\(d=3\)ではゲージ変換で候補法線が動かず、\(d\ge4\)では射影子が不変のまま候補法線がほぼ直交方向まで動く。定理3の数値的裏づけである。\(d=4,5\)は空間として採用したのではなく、非一意性の比較対象である。

\begin{figure}
\centering
\pandocbounded{\includegraphics[keepaspectratio,alt={法線ゲージ自由度}]{figs/normal_gauge_freedom_v1.png}}
\caption{法線ゲージ自由度}
\end{figure}

\subsubsection{8.6 実験D:
N=5のABCD部分系}\label{ux5b9fux9a13d-n5ux306eabcdux90e8ux5206ux7cfb}

\(N=5\)の8試行に対し、5個のABCD部分系（計40件）それぞれについて、部分系の6関係波の初期位相から同じ結合則で局所生成子を構成した。全40件で局所生成子はランク6・三回転平面となり、単独ABCD実験［4］と同一構造だった。

さらに、未規格化の局所生成子は、\(N\)体主生成子の対応する主部分行列と厳密に一致した（最大差
\(0\)）。端点共有隣接が部分系内で閉じているため、恒等的に一致する。したがって部分系読出しの実装は、主生成子の部分行列と局所再構成のどちらでも同一であり、差はスペクトル規格化の定数倍のみである。

\subsubsection{8.7 実験E:
多次元核の残余部分空間}\label{ux5b9fux9a13e-ux591aux6b21ux5143ux6838ux306eux6b8bux4f59ux90e8ux5206ux7a7aux9593}

\(N=6\)〜\(9\)の各8試行で、核射影子の冪等誤差は最大
\(1.11\times10^{-15}\)、生成子消滅誤差 \(\lVert KP_0\rVert\) は最大
\(1.81\times10^{-16}\) だった。核次元は全試行で \(M-2N\) に一致した。

独立に選んだ二つの正規直交核基底に対し、射影子の差は最大
\(6.66\times10^{-16}\)
であり射影子は一意である。一方、基底の第一方向どうしの角度は最小
\(23.5^\circ\)
であり、内部基底は一意でない。多次元核が「射影子と二乗量は読めるが、内部の各方向は選べない残余部分空間」であることの直接の実証である。

\begin{figure}
\centering
\pandocbounded{\includegraphics[keepaspectratio,alt={残余部分空間の成長}]{figs/kernel_dimension_growth_v1.png}}
\caption{残余部分空間の成長}
\end{figure}

\begin{center}\rule{0.5\linewidth}{0.5pt}\end{center}

\subsection{9.
本稿の主張範囲}\label{ux672cux7a3fux306eux4e3bux5f35ux7bc4ux56f2}

\subsubsection{9.1
証明した数学的帰結}\label{ux8a3cux660eux3057ux305fux6570ux5b66ux7684ux5e30ux7d50}

\begin{enumerate}
\def\labelenumi{\arabic{enumi}.}
\tightlist
\item
  頂点分解
  \(\widetilde K=\sum_k(c_ks_k^{\mathsf T}-s_kc_k^{\mathsf T})\)（定理1）。
\item
  全\(N\)・全位相でのランク上界
  \(\operatorname{rank}K\le2\min(N,\lfloor M/2\rfloor)\)（系1）。
\item
  \(3\le N\le12\)の一般位置等号（定理2、厳密有理数証人＋測度零論法）。
\item
  法線の符号を除く一意性が \(d=3\) に限ること（定理3）。
\end{enumerate}

\subsubsection{9.2
数値実験で確認した事実}\label{ux6570ux5024ux5b9fux9a13ux3067ux78baux8a8dux3057ux305fux4e8bux5b9f}

\begin{enumerate}
\def\labelenumi{\arabic{enumi}.}
\tightlist
\item
  \(N=3\)〜\(9\)、32試行×720ステップで、閉鎖・絶対値二乗和・直交性・名称置換共変性が許容差
  \(10^{-10}\) 内で維持された。
\item
  一般位置ランク則が\(N=3\)〜\(12\)の2560試行で例外なく成立した。
\item
  二値位相 \(\{0,0.5\}\) は \(M=2N\) の臨界構成 \(N=5\)
  でのみランクを低下させた。
\item
  \(N=5\)の全ABCD部分系が三回転平面構造を持ち、局所生成子は主生成子の主部分行列と厳密に一致した。
\item
  \(N\ge6\)の多次元核で、射影子の一意性と内部基底の非一意性が同時に成立した。
\end{enumerate}

\subsubsection{9.3
本稿が採用する空間方向解釈}\label{ux672cux7a3fux304cux63a1ux7528ux3059ux308bux7a7aux9593ux65b9ux5411ux89e3ux91c8}

\begin{enumerate}
\def\labelenumi{\arabic{enumi}.}
\tightlist
\item
  区別可能な三回転平面までをXYZ三方向の読出しとする（［4］の継承）。
\item
  一意性条件を超える回転平面を内部位相モードとする。
\item
  多次元核を、一意な基底を持たない残余部分空間とする。
\end{enumerate}

\subsubsection{9.4
本稿で固定しないもの}\label{ux672cux7a3fux3067ux56faux5b9aux3057ux306aux3044ux3082ux306e}

\begin{enumerate}
\def\labelenumi{\arabic{enumi}.}
\tightlist
\item
  全\(N\)に対する一般位置等号（新仮説として明示）。
\item
  残余部分空間と内部位相モードが対応する物理軸。
\item
  異なる部分系の三方向構造が動力学によって同じXYZへ位相固定される機構。
\item
  二体関係以外の高次関係を追加した場合の三方向上限の存否。
\item
  位相差正弦生成子そのものの公理からの導出。
\item
  逐次再構成生成子（各ステップで位相から生成子を再構成する動力学）の性質。本稿の対象は固定生成子である。
\end{enumerate}

\begin{center}\rule{0.5\linewidth}{0.5pt}\end{center}

\subsection{10. 再現可能性}\label{ux518dux73feux53efux80fdux6027}

本稿の数値実験は、次の実プログラムと出力に基づく。

\subsubsection{実行プログラム}\label{ux5b9fux884cux30d7ux30edux30b0ux30e9ux30e0}

\begin{itemize}
\tightlist
\item
  \texttt{次元の生成構造/対照実験\_一角度円周位相調和読出し\_v1/run\_nbody\_rank\_saturation\_preliminary\_v1.py}（実験A〜E）
\item
  \texttt{次元の生成構造/対照実験\_一角度円周位相調和読出し\_v1/run\_nbody\_exact\_rational\_witness\_v1.py}（厳密有理数証人）
\item
  \texttt{次元の生成構造/対照実験\_一角度円周位相調和読出し\_v1/make\_nbody\_rank\_saturation\_figures\_v1.py}（図生成）
\end{itemize}

\subsubsection{主要出力}\label{ux4e3bux8981ux51faux529b}

\begin{itemize}
\tightlist
\item
  \texttt{nbody\_rank\_saturation\_preliminary\_result\_v1.json}
\item
  \texttt{experiment\_a\_trials\_v1.csv} /
  \texttt{experiment\_a\_summary\_v1.csv}
\item
  \texttt{experiment\_b\_rank\_law\_v1.csv}
\item
  \texttt{experiment\_c\_normal\_uniqueness\_v1.csv}
\item
  \texttt{experiment\_d\_subsystems\_v1.csv}
\item
  \texttt{experiment\_e\_kernel\_v1.csv}
\item
  \texttt{nbody\_exact\_rational\_witness\_v1.json}
\item
  \texttt{nbody\_rank\_saturation\_preliminary\_report\_v1.md}
\end{itemize}

これらは次のディレクトリに保存されている。

\texttt{次元の生成構造/figs/}

\begin{center}\rule{0.5\linewidth}{0.5pt}\end{center}

\subsection{11. 最終結論}\label{ux6700ux7d42ux7d50ux8ad6}

\(N\)体の完全二体関係波系を、背景空間の軸を先に置かず、公理16の完全二体関係波と二乗閉鎖を基礎とし、位相差正弦型生成子を保存的な作業仮説として構成した［1］。

内部関係波数は \(M_N=\binom{N}{2}\)
で二乗的に増える。しかし生成子は接続行列の行ごとのランク2反対称行列の和へ厳密に分解され、

\[
\operatorname{rank}K\le2\min\!\left(N,\left\lfloor\frac{M_N}{2}\right\rfloor\right)
\]

が全\(N\)で成立する。\(3\le N\le12\)では、厳密有理数の証人配置と実解析関数の零点集合の測度零性により、一般位置でこの上界が等号となることを証明した。回転モードは高々線形にしか増えない。

二関係面から法線を読む方向読出しでは、法線候補空間は \(d-2\)
次元であり、符号を除く一意性は \(d=3\)
に限る。区別可能な回転平面が三を超えて増えても、超える平面は空間方向として採用できず、内部位相モードとして残る。\(N\ge6\)で再出現する多次元核は、射影子は一意だが内部基底を選べない残余部分空間である。

したがって、本稿の結論は次である。

\[
\boxed{
\begin{aligned}
&\text{内部関係波は }O(N^2)\text{ で増える}\\
&\text{内部回転モードは }O(N)\text{ を超えない}\\
&\text{現在の二体位相関係読出し規則のもとで、一意に読める空間方向は三方向で飽和する}
\end{aligned}
}
\]

残余部分空間と内部位相モードが物理的にどの軸へ対応するかは、本稿の対象外である。

\begin{center}\rule{0.5\linewidth}{0.5pt}\end{center}

\subsection{参考文献}\label{ux53c2ux8003ux6587ux732e}

\subsubsection{自己引用}\label{ux81eaux5df1ux5f15ux7528}

\begin{enumerate}
\def\labelenumi{\arabic{enumi}.}
\tightlist
\item
  木原 範昭「無名等振幅複合波モデル基本公理系 v5」Zenodo, 2026. Version
  DOI:
  \href{https://doi.org/10.5281/zenodo.21465429}{10.5281/zenodo.21465429},
  Concept DOI:
  \href{https://doi.org/10.5281/zenodo.21315735}{10.5281/zenodo.21315735}.
\item
  木原
  範昭「AB二体閉鎖位相系における調和読出しとc=1面積スイープ予備実験総括
  v4」Zenodo, 2026. Version DOI:
  \href{https://doi.org/10.5281/zenodo.21374317}{10.5281/zenodo.21374317},
  Concept DOI:
  \href{https://doi.org/10.5281/zenodo.21318696}{10.5281/zenodo.21318696}.
\item
  木原
  範昭「AB二体閉鎖位相系における未来位相位置加速度写像と調和閉鎖による逆二乗則」Zenodo,
  2026. Version DOI:
  \href{https://doi.org/10.5281/zenodo.21441082}{10.5281/zenodo.21441082},
  Concept DOI:
  \href{https://doi.org/10.5281/zenodo.21441081}{10.5281/zenodo.21441081}.
\item
  木原
  範昭「完全二体関係波から読み出されるXYZ三方向――AB・ABC・ABCD閉鎖系における内部関係方向の増加と空間方向読出しの飽和」Zenodo,
  2026. Version DOI:
  \href{https://doi.org/10.5281/zenodo.21454790}{10.5281/zenodo.21454790},
  Concept DOI:
  \href{https://doi.org/10.5281/zenodo.21454789}{10.5281/zenodo.21454789}.
\end{enumerate}

\subsubsection{外部文献}\label{ux5916ux90e8ux6587ux732e}

\begin{enumerate}
\def\labelenumi{\arabic{enumi}.}
\setcounter{enumi}{4}
\tightlist
\item
  Hassler Whitney, ``Congruent Graphs and the Connectivity of Graphs,''
  \emph{American Journal of Mathematics}, 54(1), 150--168, 1932. DOI:
  \href{https://doi.org/10.2307/2371086}{10.2307/2371086}.
\item
  Milan Vujivčić, Fedor Herbut, and Gradimir Vujivčić, ``Canonical Form
  for Matrices Under Unitary Congruence Transformations. I:
  Conjugate-Normal Matrices,'' \emph{SIAM Journal on Applied
  Mathematics}, 23(2), 225--238, 1972. DOI:
  \href{https://doi.org/10.1137/0123025}{10.1137/0123025}.
\item
  Fasma Diele, Luciano Lopez, and R. Peluso, ``The Cayley Transform in
  the Numerical Solution of Unitary Differential Systems,''
  \emph{Advances in Computational Mathematics}, 8(4), 317--334, 1998.
  DOI:
  \href{https://doi.org/10.1023/A:1018908700358}{10.1023/A:1018908700358}.
\item
  Cyriel van Nuffelen, ``On the incidence matrix of a graph,''
  \emph{IEEE Transactions on Circuits and Systems}, 23(9), 572, 1976.
  DOI:
  \href{https://doi.org/10.1109/TCS.1976.1084251}{10.1109/TCS.1976.1084251}.
\item
  Boris S. Mityagin, ``The Zero Set of a Real Analytic Function,''
  \emph{Mathematical Notes}, 107(3), 529--530, 2020. DOI:
  \href{https://doi.org/10.1134/S0001434620030189}{10.1134/S0001434620030189}.
\end{enumerate}

\end{document}
